\chapter{Vaginal Candidiasis}

\section*{Definition and Overview}
Vaginal candidiasis, commonly called vaginal thrush, is an infection of the vagina and the surrounding vulval area caused by \textit{Candida} yeasts, most often \textit{Candida albicans}. These yeasts live naturally in the vagina of many healthy women, but when they multiply excessively they provoke inflammation, itching, and discharge. It is a very common and generally harmless condition, although it can be uncomfortable, distressing, and frustrating when it recurs. It is not a sexually transmitted infection, although sexual activity can sometimes trigger or worsen symptoms. The term also covers related presentations such as recurrent vulvovaginal candidiasis, defined as four or more episodes in a year, which deserves a structured management approach.

\section*{Epidemiology}
Vaginal thrush is extremely common: roughly three in four women will experience at least one episode in their lifetime, and a substantial number will have recurrent episodes. It is most common in women of reproductive age, with a peak in the third decade of life. Certain situations increase the risk, including pregnancy, the use of broad-spectrum antibiotics, poorly controlled diabetes, the combined contraceptive pill, and immune suppression from illness or medication. After the menopause, symptoms may persist or reappear in women using hormone replacement therapy. Recurrence is more common in women with diabetes, in those who use antibiotics frequently, and in a minority of women with no identifiable trigger in whom the mechanisms of recurrence remain poorly understood.

\section*{Aetiology and Pathophysiology}
\textit{Candida} is part of the normal vaginal flora in many women, kept in balance by protective bacteria, mainly lactobacilli, and by the local immune system. Thrush develops when this balance is disturbed and the yeast overgrows:
\begin{itemize}
    \item \textbf{Antibiotics:} kill protective lactobacilli, allowing yeasts to flourish
    \item \textbf{Hormonal changes:} pregnancy, the menstrual cycle, and the combined pill raise oestrogen levels, which encourage yeast growth
    \item \textbf{High blood sugar:} diabetes provides glucose that feeds the yeast and impairs immune responses
    \item \textbf{Immune suppression:} from illness, chemotherapy, or corticosteroid treatment
    \item \textbf{Local irritation:} tight synthetic underwear, perfumed products, or prolonged moisture, which damage the skin barrier
\end{itemize}
The growing yeast attaches to the vaginal lining and triggers an inflammatory response, causing redness, swelling, itching, and the typical discharge. In otherwise healthy women the infection does not usually spread beyond the surface, but in severely immunocompromised women it can, rarely, progress to invasive disease.

\section*{Clinical Features}
Symptoms of vaginal thrush may include:
\begin{itemize}
    \item Intense itching and soreness of the vulva and vagina
    \item A thick, white, cottage-cheese-like discharge that is usually odourless
    \item Redness, swelling, and sometimes cracking or fissuring of the vulval skin
    \item Burning, especially when passing urine or during intercourse
    \item In some women, a thinner, watery discharge or no discharge at all
    \item Worsening of symptoms in the week before a period
    \item Swelling of the labia and a rash spreading to the inner thighs in severe cases
\end{itemize}

\section*{Differential Diagnosis}
Other conditions can cause similar symptoms and should be distinguished, especially because they need different treatment:
\begin{itemize}
    \item Bacterial vaginosis, which causes a thin grey discharge and a fishy odour
    \item Trichomoniasis, a sexually transmitted infection with a frothy, often offensive discharge
    \item Chlamydia and gonorrhoea, which can cause discharge, bleeding, and pelvic pain
    \item Irritant or allergic contact dermatitis from soaps, washes, and hygiene products
    \item Lichen sclerosus or eczema of the vulva, which cause itching and skin changes
    \item Atrophic vaginitis after the menopause, with dryness and soreness rather than discharge
\end{itemize}

\section*{When to Seek Medical Care}
Many women can treat a first or typical episode with over-the-counter antifungal products. See a doctor if:
\begin{itemize}
    \item You have never had thrush before and are unsure of the cause
    \item Symptoms are severe, do not improve with treatment, or keep returning
    \item You are pregnant or breastfeeding
    \item You have diabetes, a weakened immune system, or take immune-suppressing medicines
    \item There is an unusual discharge, a fishy smell, bleeding, or pain in the lower abdomen
    \item You have had multiple sexual partners or a partner with symptoms
\end{itemize}
\textbf{Call 000 (or local emergency services) for:}
\begin{itemize}
    \item High fever with pelvic pain
    \item Heavy vaginal bleeding or fainting
    \item Severe, sudden pelvic or abdominal pain
\end{itemize}

\section*{Diagnosis and Investigations}
\begin{itemize}
    \item \textbf{History:} careful review of symptoms, triggers such as recent antibiotics, and risk factors including diabetes and pregnancy
    \item \textbf{Examination:} inspection of the vulva and vagina to assess redness, swelling, and discharge
    \item \textbf{Vaginal swab:} for microscopy and culture, especially when the diagnosis is uncertain, treatment has failed, or infection recurs
    \item \textbf{pH testing:} of the discharge to help distinguish thrush, which is usually acidic, from bacterial vaginosis
    \item \textbf{Blood tests:} for diabetes or other underlying conditions in recurrent cases
\end{itemize}

\section*{Management}
Treatment aims to relieve symptoms and restore the balance of vaginal organisms:
\begin{itemize}
    \item \textbf{Antifungal creams:} applied to the vulva for itching and soreness
    \item \textbf{Vaginal pessaries:} antifungal tablets inserted into the vagina, usually used for three to six nights
    \item \textbf{Oral antifungal tablets:} a single dose of an oral medicine such as fluconazole is often the simplest option
    \item \textbf{Combined products:} an antifungal cream combined with a mild corticosteroid for severe inflammation
    \item \textbf{Treatment during pregnancy:} creams and pessaries are preferred, while oral tablets are generally avoided
    \item \textbf{Recurrent thrush:} longer courses of treatment followed by intermittent or weekly preventive antifungal treatment for several months, together with treatment of any underlying cause
    \item \textbf{Partner treatment:} usually not needed, as thrush is not sexually transmitted
\end{itemize}

\section*{Complications}
\begin{itemize}
    \item Recurrent or chronic thrush causing distress and disruption to daily life, sleep, and sexual relationships
    \item Cracking and fissuring of the vulval skin, with secondary bacterial infection
    \item Severe inflammation and discomfort, particularly in diabetes or pregnancy
    \item In women with a very weak immune system, rare progression of \textit{Candida} beyond the vagina
    \item Side effects of treatment, such as local irritation from creams or rare liver effects with oral antifungals
    \item Misdiagnosis and undertreatment of other conditions such as trichomoniasis or bacterial vaginosis
\end{itemize}

\section*{Prognosis and Outcomes}
With the correct treatment, symptoms of vaginal thrush usually settle within a few days, although a minority of women develop recurrent episodes that require longer-term management. Recurrent thrush can be controlled in most women with a maintenance regimen and by addressing triggers such as antibiotic use and blood sugar control. The condition does not affect fertility or general health, and serious complications are very rare in otherwise healthy women.

\section*{Prevention and Self-Care}
\begin{itemize}
    \item Wear cotton underwear and loose clothing, and avoid tight synthetic garments
    \item Avoid perfumed soaps, bubble baths, douching, and deodorant products in the genital area
    \item Wash with plain water and dry the area gently
    \item Change out of wet swimwear or sweaty clothing promptly
    \item Wipe from front to back after toileting
    \item Manage blood sugar carefully if you have diabetes
    \item Use antibiotics only when they are prescribed and needed
    \item For recurrent thrush, talk to a doctor about a preventive plan rather than self-treating repeatedly
\end{itemize}

\section*{Sources and Further Reading}
The following sources provide authoritative, up-to-date information on vaginal candidiasis and its management.
\begin{itemize}
    \item NHS. \textit{Thrush in men and women.} https://www.nhs.uk/conditions/thrush/
    \item Mayo Clinic. \textit{Vaginal yeast infection (vaginal candidiasis).} https://www.mayoclinic.org/diseases-conditions/yeast-infection/symptoms-causes/syc-20378999
    \item Centers for Disease Control and Prevention. \textit{Candidiasis: vaginal candidiasis and recurrent vulvovaginal candidiasis.} https://www.cdc.gov/fungal/diseases/candidiasis/
    \item MSD Manual. \textit{Candidiasis of the vulva and vagina.} https://www.msdmanuals.com/
    \item MedlinePlus. \textit{Vaginal yeast infection.} https://medlineplus.gov/yeastinfections.html
\end{itemize}
